% II. Проектиране на системата.
\section{Проектиране на системата}
\label{sec:design}

\subsection{Структура на интерфейса}

Интерфейсът е конвейер от четири анализиращи слоя и един синтезиращ. Входът е низ на
естествен език, изходът е формална заявка заедно с перифраза на разбрания въпрос
(Фиг.~\ref{fig:pipeline}).

Изискванията са пет. Всеки слой приема и връща данни по описан договор и може да бъде
заменен без промяна в останалите. Нееднозначността се пренася нагоре, вместо да се
разрешава преждевременно с евристика на ниско ниво. Всяко тълкуване е проследимо назад до
думите, които са го породили. Заявката не се изпълнява преди потвърждение. И накрая, цялото
знание за езика и за предметната област живее във файлове с данни, тоест смяната на
предметната област е смяна на данни, а не на изходен текст.

\begin{figure}[H]
\centering
\begin{tikzpicture}[node distance=0.7cm]
    \node[block] (q) {Въпрос на естествен език};
    \node[note, left=1.6cm of q] (lex) {\code{lexicon\_en.txt}\\лексикон и парадигми};
    \node[block, below=of q] (morph) {Морфологичен анализ\\краен преобразувател};
    \node[block, below=of morph] (parse) {Синтактичен анализ\\Earley, гора от изводи};
    \node[note, left=1.6cm of parse] (gram) {\code{grammar\_en.txt}\\граматика};
    \node[block, below=of parse] (sem) {Семантично изображение\\типова проверка};
    \node[note, left=1.6cm of sem] (schema) {\code{schema.txt}\\модел на знанието};
    \node[block, below=of sem] (lf) {Логическа форма};
    \node[accent, below left=1.0cm and 0.4cm of lf] (sql) {Генератор на SQL};
    \node[accent, below right=1.0cm and 0.4cm of lf] (sparql) {Генератор на SPARQL};
    \node[accent, right=2.0cm of lf] (synth) {Синтез на изречение\\перифраза};
    \node[danger, right=1.6cm of synth] (user) {Потвърждение\\от потребителя};

    \draw[line] (lex) -- (morph);
    \draw[line] (q) -- (morph);
    \draw[line] (morph) -- node[right, font=\scriptsize] {разчитания за всяка дума} (parse);
    \draw[line] (gram) -- (parse);
    \draw[line] (parse) -- node[right, font=\scriptsize] {гора от изводи} (sem);
    \draw[line] (schema) -- (sem);
    \draw[line] (sem) -- (lf);
    \draw[line] (lf) -- (sql);
    \draw[line] (lf) -- (sparql);
    \draw[line] (lf) -- (synth);
    \draw[line] (synth) -- (user);
\end{tikzpicture}
\caption{Структура на конвейера и поток на данните между слоевете}
\label{fig:pipeline}
\end{figure}

\subsection{Морфологичен анализ}

Слоят връща за всяка словоформа множество от разчитания, всяко с лексема, част на речта и
граматични признаци. Множество, а не единично разчитане, защото омонимията е обичайна и се
разрешава едва по-нагоре: \term{films} се анализира и като множествено число на
съществителното, и като лична форма на глагола.

Преобразувателят е композиция на две релации: лексиконът, тоест основи с лексема, част на
речта и парадигма, и афиксната релация на парадигмата, тоест форми във вида „изтрий $n$
знака от основата, после добави този суфикс“ заедно с признаците, които афиксът внася.
Композицията се извършва еднократно при зареждане, дава детерминиран ацикличен
преобразувател, в който всеки път изписва една словоформа, и се пази и в обратната посока.

Две решения заслужават коментар. Изтриването в афикса се брои в знаци, а не в байтове:
английският лексикон не би забелязал разликата, българският би, защото кирилската буква
заема два байта в UTF-8. И редуването живее в данните: \code{\textasciitilde 1+ies} казва
„махни един знак, добави \code{ies}“, а \code{\textasciitilde 6+people} превръща
\code{person} в \code{people}. Неправилната форма е ред в лексикона, а не нов C++.

\subsection{Синтактичен анализ}

Слоят приема поредица от множества с разчитания и връща гора от изводи. Терминал на
граматиката не е дума, а част на речта с ограничение върху признаците, затова граматиката
е независима от лексикона.

Елементите на таблицата на Earley се обединяват по ключ „продукция, позиция на точката,
начало“ за всяка крайна позиция, а всеки елемент носи списък от връзки. Една връзка е един
начин да се стигне до елемента: предшественикът с точка една позиция вляво и възелът, който
запълва символа между тях. Нееднозначното изречение затова умножава връзки, а не елементи,
и точно това прави гората уплътнена. Възелът се идентифицира с тройката „символ, начало,
край“, така че еднаквите подфрази се споделят. Броят на изводите се пресмята върху гората
чрез динамично оптимиране, а не чрез изброяване; граматиката няма празни продукции и няма
цикъл от единични продукции, откъдето гората е ациклична и броят е краен.

Нееднозначността по прикрепване живее в точно три продукции: \code{NOM -> NOM PP},
\code{VP -> VP PP} и \code{VPPASS -> VPPASS PP}. Нищо друго не поражда две дървета за един
и същ низ, така че отчетената нееднозначност е нееднозначност по прикрепване, а не
артефакт от излишна продукция.

\subsection{Семантичен анализ}

Изображението е композиционно: всяка продукция носи име на семантично правило, а правилото
изгражда значението на групата от значенията на съставките ѝ. Тук разчитанията отпадат.
Всеки атом се проверява типово спрямо модела, а предлогът се превръща в релация само ако
моделът съдържа \term{рамка} за него, тоест тройката „предлог, тип на аргумента, тип на
носителя“ заедно с релацията и реда на аргументите. Затова \term{in 2010} се прикрепя към
филм, а \term{in Inception} към лице: двата се различават по типа на именната група, а не
по правило за думата \term{in}. Разчитане, което не се изобразява върху нищо в модела, се
отхвърля с посочена причина, и това отхвърляне, а не ранна евристика, разрешава
нееднозначността.

Когато оцелее повече от едно разчитане, изборът се прави по предпочитание за близко
прикрепване: сумата от разстоянията между всяка предложна група и главата, към която се
прикрепя, като по-малката сума печели. Предпочитанието се прилага след семантичния филтър,
никога преди него. Разчитанията, чиито логически форми съвпадат след канонично
преименуване на променливите, се събират в едно, така че показаният брой е брой на
значенията, а не на дърветата.

\subsection{Генериране на формална заявка}

Генераторът работи изцяло върху проверената логическа форма и не съдържа евристики върху
текста, затова една и съща семантика се насочва към две цели без промяна в предходните
слоеве. Съответствието е пряко: едноместният атом става \code{FROM} с псевдоним и тройка с
\code{a}; двуместният става съединение по таблица или равенство върху колона, срещу една
тройка; проекцията става \code{SELECT DISTINCT} и в двата. Броенето е
\code{COUNT(DISTINCT ...)}, въпросът с отговор да или не е \code{EXISTS} и \code{ASK},
отрицанието е \code{NOT EXISTS} и \code{FILTER NOT EXISTS}, а всеобщността е релационно
деление с две вложени отрицания.

\subsection{Синтез на изречение}
\label{sub:synth}

Синтезът произвежда изречение от логическата форма и го показва като перифраза на разбрания
въпрос. Ролята му е контролна: при грешно избрано разчитане потребителят вижда това преди
изпълнението. Перифразата се изгражда от формата, не от въпроса; иначе щеше да се съгласи с
потребителя по построение.
